% Generated draft: requires mathematical and source review.
\documentclass[11pt]{article}
\usepackage[margin=25mm]{geometry}
\usepackage{fontspec}
\setmainfont{Latin Modern Roman}
\usepackage{amsmath,amssymb,mathtools,unicode-math}
\setmathfont{Latin Modern Math}
\usepackage{enumitem,needspace,xurl,hyperref}
\setlength{\parindent}{0pt}
\setlength{\parskip}{5pt}
\setlength{\emergencystretch}{3em}
\newcommand{\N}{\mathbb{N}}
\newcommand{\Z}{\mathbb{Z}}
\newcommand{\Q}{\mathbb{Q}}
\newcommand{\R}{\mathbb{R}}
\newcommand{\C}{\mathbb{C}}
\newcommand{\F}{\mathbb{F}}
\newcommand{\A}{\mathbb{A}}
\newcommand{\PP}{\mathbb{P}}
\newcommand{\E}{\mathbb{E}}
\newcommand{\eps}{\varepsilon}
\newcommand{\dd}{\,\mathrm{d}}
\begin{document}
\section*{UnsolvedMath ID: 30003265}
\textbf{Source problem code:} catalogue entry
\subsection{Definitions and mathematical statement}
Let \(\mathcal O=\Z[i]\) be the ring of Gaussian integers, and let \(\Gamma=\operatorname{SL}_2(\mathcal O)\). For a positive rational prime \(p\equiv5\pmod8\), write
\[
(p)=\mathfrak p\,\overline{\mathfrak p},
\]
where \(\mathfrak p\) is a prime ideal of \(\mathcal O\) and \(\overline{\mathfrak p}\) is its complex conjugate. Here the factorization is an equality of ideals; it makes explicit the ideal interpretation of the source's notation \(p=\mathfrak p\overline{\mathfrak p}\). Define
\[
\Gamma_0(\mathfrak p)=
\left\{
\begin{pmatrix}a&b\\c&d\end{pmatrix}
\in\operatorname{SL}_2(\mathcal O):c\in\mathfrak p
\right\}.
\]
Let \(\mathbb H^3=\{(z,t):z\in\C,\ t\in\R_{>0}\}\) be hyperbolic three-space in its upper half-space model, with the standard action of \(\operatorname{SL}_2(\C)\) through \(\operatorname{PSL}_2(\C)\), and put
\[
X_{\mathfrak p}=\Gamma_0(\mathfrak p)\backslash\mathbb H^3.
\]
As written, \(H_1(X_{\mathfrak p};\Z)\) denotes ordinary singular homology of the orbit space with integer coefficients.

The supplied question asks whether there exist such a prime \(p\), a prime ideal \(\mathfrak p\) above \(p\), and a nonzero class
\[
\alpha\in H_1(X_{\mathfrak p};\Z)
\]
that is \emph{cuspidal}, or whether no such class exists for any prime \(p\equiv5\pmod8\) and either prime ideal above it.

The term \emph{cuspidal class} is not defined in the supplied record. In particular, the record does not specify an integral cuspidal subgroup or quotient, whether torsion classes count, or whether cuspidality is tested after extending coefficients to \(\Q\) or \(\C\). Thus the existential question and its universal negative are preserved here, but a fully determined mathematical assertion requires that convention to be supplied. The source also does not clarify whether its computations concern the displayed orbit-space homology or group or orbifold homology; these interpretations must not be silently interchanged with integer coefficients.
\subsection{Short English statement}
The question asks whether a Gaussian Bianchi congruence quotient of prime ideal level above a rational prime \(p\equiv5\pmod8\) has a nonzero cuspidal class in its first integral homology, or whether this never occurs. The meaning of integral cuspidality and the intended homology interpretation require clarification. The supplied record labels the problem open; this draft does not independently verify that status.
\subsection{Sources}
\begin{itemize}
\item Supplied local UnsolvedMath registry record, permanent ID \(30003265\), title: \emph{Cuspidal Homology for Gaussian Bianchi Congruence Groups}; exact source problem code: \texttt{catalogue entry}. Collection: Oberwolfach Reports Open Problems. No individual proposer is identified; the proposed year is recorded as 2016.
\item Supplied report citation: \emph{Mini-Workshop: Computations in the Cohomology of Arithmetic Groups} (2017), Oberwolfach Reports, pp.~2941--2973, DOI \texttt{10.4171/owr/2016/52}. \url{https://doi.org/10.4171/owr/2016/52}. The supplied background describes these reports as documentation of mathematical workshops and their open-problem discussions.
\item Supplied literature source: \emph{OWR 52 preliminary report}. \url{https://orbilu.uni.lu/bitstream/10993/29499/1/preliminary_OWR_2016_52.pdf}.
\item Literature assessment dated exactly \texttt{2026-08-21}: the supplied curation labels the problem open and reports that the computation found no cuspidal class for \(p<16000\); it also reports that no proof of nonexistence for all eligible primes, and no example beyond the search range, was found. This is supplied dated curation, not an independently verified literature assessment. The supplied background explicitly requests independent verification of time-sensitive status and cited claims.
\item The registry labels the statement \texttt{corrected\_verified} and says that mathematical content was preserved while notation, typography, or omitted definitions were repaired. These are supplied editorial assertions, not independent mathematical verification of this draft.
\end{itemize}
\end{document}
